% Abstract v2 - the equivalence/parsimony thesis. Replaces the v1 "mass-aware
% / +61% vs CONGA" abstract.
\begin{abstract}
We ask how far classical physics goes in network load balancing on its own.
Modelling a packet as a particle descending a potential field, we reduce a
routing rule to two mechanisms: gradient descent toward less latency and
congestion, and a Newton's-third-law reaction in which a dropped packet leaves
a decaying ``scar'' steering its successors away. An ablation of a richer
eight-layer model (effective mass, thermostat, temporal ripple, buoyancy)
shows that only these two terms earn their keep; the rest are inert or mildly
harmful. We then test, as a statistical \emph{equivalence} claim (TOST,
margin $2$ percentage points), whether this two-term core matches purpose-built
load balancers. Across fifteen topologies in a flow-level simulator, the core
is statistically equivalent to DRILL on twelve, and---once CONGA is given a
fair path budget---equivalent to CONGA on the regular topologies, reaching the
same operating point that CONGA attains only after enumerating and scoring
$16$--$32$ candidate paths per flow. The result is equivalence with parsimony,
not dominance: the physics needs no cleverness beyond two classical forces. We
further show that a purely topological quantity, the tube/sp ratio (room to
manoeuvre among near-shortest paths, computable before deployment), predicts
where the equivalence holds and where it breaks (Pearson $r = 0.78$): the two
real WAN backbones, with low tube/sp, are exactly the cases where engineered
multi-path routing retains a genuine edge. We report negative results and a
self-audit of an earlier inflated finding in full, as part of the contribution.
\end{abstract}
